\documentclass{article}

\usepackage{amsmath,amssymb}
\usepackage{algorithm,algorithmic}
\usepackage{graphicx}
\usepackage{booktabs}
\usepackage{multirow}
\usepackage{xcolor}
\usepackage{url}

\newcommand{\todo}[1]{{\color{red}[#1]}}
\newcommand{\R}{\mathbb{R}}

\title{Microstructure-Aware Conservative Execution via Reinforcement Learning with Group Relative Policy Optimization}

\author{
    Anonymous Authors \\
    Institution Name Withheld for Review
}

\begin{document}

\maketitle

\begin{abstract}
Optimal trade execution requires balancing market impact, opportunity cost, and microstructure effects. While reinforcement learning (RL) offers a promising framework, traditional actor-critic methods suffer from instability in offline settings where value function approximation is unreliable. We propose \textbf{MACE-RL} (Microstructure-Aware Conservative Execution via RL), a novel offline RL framework that combines a normalizing-flow action manifold, microstructure-aware state encoding, and DeepSeek-style Group Relative Policy Optimization (GRPO). Our method learns execution policies directly from historical limit-order-book (LOB) data without a value network, using group sampling to estimate advantages relative to alternative actions at each state. The policy consists of: (1) a microstructure encoder capturing spread, imbalance, depth, and volatility; (2) a conditional normalizing-flow defining a feasible action manifold learned from historical execution patterns; (3) a residual adaptation module responding to liquidity shocks while respecting the manifold. We implement a realistic execution environment using the FI-2010 LOB dataset with implementation shortfall rewards and market impact. Experiments demonstrate that MACE-RL outperforms traditional execution benchmarks and standard offline RL methods, providing stable, microstructure-aware execution strategies. Our contributions include the first application of GRPO to optimal execution, the integration of normalizing flows with residual adaptation for action feasibility, and a theoretically grounded approach to offline policy optimization without critic networks.
\end{abstract}

\section{Introduction}
Optimal execution—the problem of liquidating or acquiring a large asset position while minimizing transaction costs—remains a fundamental challenge in quantitative finance. The classical Almgren-Chriss model \citep{almgren2001optimal} provides an analytical framework but assumes linear market impact and ignores microstructure effects. In practice, execution occurs in discrete time with observable limit-order-book (LOB) dynamics, including spread, depth, imbalance, and volatility \citep{gould2013limit}. These microstructure features contain valuable signals for adaptive execution strategies.

Reinforcement learning (RL) offers a natural framework for learning adaptive execution policies from historical data \citep{neuneier1998optimal, nevmyvaka2006reinforcement}. However, traditional actor-critic methods face significant challenges in offline settings: value function approximation becomes unreliable due to distributional shift \citep{fujimoto2019off}, leading to overestimation and training instability. Recent offline RL methods like Conservative Q-Learning (CQL) \citep{kumar2020conservative} and Implicit Q-Learning (IQL) \citep{kostrikov2022offline} introduce conservatism but still rely on value networks, which may propagate errors in financial domains with sparse, noisy rewards.

We address these limitations with a \textit{critic-free} approach based on DeepSeek's Group Relative Policy Optimization (GRPO) \citep{deepseek2024grpo}. GRPO eliminates the value network entirely, instead sampling multiple actions per state and evaluating their relative advantages directly. This paradigm shift aligns well with execution problems, where evaluating alternative execution rates at a given market state is computationally feasible. Our method, MACE-RL, makes the following contributions:

\begin{itemize}
    \item \textbf{First application of DeepSeek-style GRPO to optimal execution}, providing stable offline policy optimization without value function approximation errors.
    \item \textbf{Integration of normalizing-flow action manifolds with residual adaptation}, enabling feasible action spaces learned from historical data while allowing adaptive responses to liquidity shocks.
    \item \textbf{Microstructure-aware state encoding} that captures spread, depth, imbalance, volatility, and liquidity indicators from LOB data.
    \item \textbf{Theoretical motivation} for group-relative advantages as a form of microstructure-driven robustness, comparing execution strategies under identical market conditions.
\end{itemize}

\section{Related Work}

\subsection{Optimal Execution}
The optimal execution literature originates from \citet{almgren2001optimal}, who model market impact as linear and derive continuous-time solutions. Time-weighted average price (TWAP) and volume-weighted average price (VWAP) are heuristic benchmarks \citep{bergault2021multi}. More recent approaches incorporate LOB dynamics \citep{cartea2015algorithmic} and machine learning \citep{lim2020reinforcement}. Our work extends this line by learning microstructure-aware policies directly from historical LOB data.

\subsection{Market Microstructure Learning}
LOB modeling has evolved from analytical models \citep{cont2010stochastic} to machine learning approaches \citep{tsantekidis2020deep}. Deep learning architectures capture temporal dependencies in LOB data \citep{zhang2021temporal}. We build on this literature by using LOB features as RL state representations.

\subsection{Offline Reinforcement Learning}
Offline RL addresses the challenge of learning from fixed datasets without environment interaction. Conservative methods like CQL \citep{kumar2020conservative}, BCQ \citep{fujimoto2019off}, and IQL \citep{kostrikov2022offline} constrain policy updates to prevent extrapolation errors. However, they retain critic networks. GRPO \citep{deepseek2024grpo} represents a distinct approach that eliminates the critic entirely, which we adapt to the execution domain.

\subsection{Normalizing Flows in RL}
Normalizing flows \citep{rezende2015variational} provide invertible transformations for density estimation. In RL, flows have been used for action representation \citep{ward2020improving}, exploration \citep{mazoure2020deep}, and offline policy constraints \citep{siegel2020keep}. Our use of conditional flows to define feasible action manifolds follows this line but specifically targets execution constraints.

\subsection{Group Policy Optimization}
The GRPO algorithm \citep{deepseek2024grpo} samples multiple actions per state and computes relative advantages within each group. This approach shares conceptual roots with evolutionary strategies \citep{salimans2017evolution} and estimation-of-distribution algorithms but is tailored for policy gradient methods. To our knowledge, ours is the first application to financial problems.

\subsection{Execution Simulators}
Realistic execution simulators require modeling market impact, LOB dynamics, and transaction costs \citep{bacry2015hawkes, donier2015won}. We build a simulator using the FI-2010 dataset \citep{nystrom2014high}, which provides realistic LOB trajectories for evaluation.

\section{Method}
\label{sec:method}

Our method, MACE-RL, consists of five components: (1) microstructure state representation, (2) normalizing-flow action manifold, (3) residual execution module, (4) execution environment and reward, and (5) DeepSeek GRPO optimization. We describe each in detail.

\subsection{Microstructure State Representation}
\label{subsec:state}

Let $s_t \in \R^d$ denote the microstructure state at time $t$. We extract features from the LOB snapshot at time $t$:

\begin{itemize}
    \item \textbf{Spread}: Difference between best ask and bid prices.
    \item \textbf{Depth}: Available volume at top $k$ price levels.
    \item \textbf{Imbalance}: Ratio of bid to ask volumes.
    \item \textbf{Volatility}: Recent price movements (e.g., realized variance).
    \item \textbf{Liquidity shocks}: Abnormal volume or order flow.
    \item \textbf{Momentum}: Short-term price trends.
\end{itemize}

These features are concatenated into $s_t$. A multilayer perceptron (MLP) encoder $h_\phi(s_t)$ produces a latent representation used by subsequent modules:

\begin{equation}
    h_t = \text{MLP}_\phi(s_t) \in \R^{d_h}.
\end{equation}

\subsection{Normalizing-Flow Action Manifold}
\label{subsec:flow}

The action space consists of execution rates $a_t \in [0,1]$, representing the fraction of remaining volume to execute at step $t$. To constrain actions to historically feasible patterns, we learn a conditional normalizing flow \citep{rezende2015variational}.

Let $z \sim \mathcal{N}(0, I)$ be a latent variable. The flow $f_\theta(z; h_t)$ maps $z$ to an action $a_t$ conditioned on the state encoding $h_t$. Using RealNVP \citep{dinh2017density} with $K$ coupling layers:

\begin{equation}
    a_t = f_\theta(z; h_t), \quad \text{with} \quad \log p(a_t | h_t) = \log \mathcal{N}(z; 0, I) - \log \left| \det \frac{\partial f_\theta}{\partial z} \right|.
\end{equation}

The flow is trained via maximum likelihood on historical execution data, learning a manifold of actions that respect historical execution patterns. During policy learning, actions are sampled by drawing $z \sim \mathcal{N}(0, I)$ and transforming through $f_\theta$.

\subsection{Residual Execution Module}
\label{subsec:residual}

To adapt to transient liquidity conditions not captured by the static flow manifold, we add a residual module $r_\psi(h_t)$ that outputs a small adjustment $\delta_t \in [-0.1, 0.1]$. The final action becomes:

\begin{equation}
    a_t = \text{clip}\left( f_\theta(z; h_t) + \lambda r_\psi(h_t), \; 0, 1 \right),
\end{equation}
where $\lambda = 0.1$ scales the residual. The residual module is a shallow MLP with tanh output activation.

The log probability of the action must account for the deterministic residual shift. Let $a_t^{\text{base}} = a_t - \lambda r_\psi(h_t)$. Then:

\begin{equation}
    \log \pi_\theta(a_t | s_t) = \log p(a_t^{\text{base}} | h_t) = \log \mathcal{N}(z; 0, I) - \log \left| \det \frac{\partial f_\theta}{\partial z} \right|,
\end{equation}
where $z = f_\theta^{-1}(a_t^{\text{base}}; h_t)$.

\subsection{Execution Environment and Reward}
\label{subsec:env}

We implement a discrete-time execution environment based on the FI-2010 LOB dataset \citep{nystrom2014high}. Each episode corresponds to executing a volume $V$ over a trajectory of LOB states $\{s_0, \dots, s_T\}$.

\textbf{State transition}: The environment progresses through the LOB trajectory. The agent's remaining volume $v_t$ updates as $v_{t+1} = v_t - a_t v_t$.

\textbf{Execution price}: Given mid-price $m_t$ (derived from LOB), the execution price includes temporary and permanent impact:
\begin{equation}
    p_t^{\text{exec}} = m_t + \eta \cdot (a_t v_t) + \psi \cdot (V - v_t),
\end{equation}
where $\eta$ and $\psi$ are impact coefficients.

\textbf{Reward}: We use implementation shortfall with risk penalty:
\begin{equation}
    r_t = \begin{cases}
        - \left( \text{cost}_t + \alpha v_t^2 \right) & \text{if terminal} \\
        - \left( \text{cost}_t + \beta v_t \right) & \text{otherwise},
    \end{cases}
\end{equation}
where $\text{cost}_t = (p_t^{\text{exec}} - m_0) \cdot (a_t v_t)$, $\alpha$ is risk aversion, and $\beta$ is a delay penalty.

\textbf{Internal state snapshot}: For GRPO group evaluation, the environment supports snapshot/restore of internal state $(t, v_t, \text{cost}_{\text{cum}})$, enabling reward evaluation for multiple actions at the same LOB state without advancing the trajectory.

\subsection{DeepSeek GRPO for Execution}
\label{subsec:grpo}

\subsubsection{No Critic Design}
GRPO eliminates the value network entirely. Advantages are computed directly from rewards of sampled actions, avoiding value function approximation errors that plague offline actor-critic methods.

\subsubsection{Group Sampling}
For each state $s_t$ in a batch, we sample $G$ actions $\{a_t^{(1)}, \dots, a_t^{(G)}\}$ from the current policy $\pi_\theta$. The first action $a_t^{(1)}$ is the rollout action (used for trajectory collection); the remaining $G-1$ are sampled anew.

\subsubsection{GRPO Advantage}
Let $r_t^{(g)}$ be the reward of action $a_t^{(g)}$ evaluated at state $s_t$ (using environment snapshot). The group-relative advantage is:
\begin{equation}
    A_t^{(g)} = \frac{r_t^{(g)} - \mu_t}{\sigma_t + \epsilon}, \quad \mu_t = \frac{1}{G} \sum_{g=1}^G r_t^{(g)}, \quad \sigma_t = \sqrt{\frac{1}{G} \sum_{g=1}^G (r_t^{(g)} - \mu_t)^2}.
\end{equation}
Normalization by group standard deviation stabilizes learning and provides scale-invariant advantages.

\subsubsection{PPO-style Loss}
Using advantages $A_t^{(g)}$ and old log probabilities $\log \pi_{\theta_{\text{old}}}(a_t^{(g)} | s_t)$, the policy loss is:
\begin{equation}
    \mathcal{L}_{\text{GRPO}} = -\frac{1}{G} \sum_{g=1}^G \min\left( \rho_t^{(g)} A_t^{(g)}, \;\text{clip}(\rho_t^{(g)}, 1-\epsilon, 1+\epsilon) A_t^{(g)} \right),
\end{equation}
where $\rho_t^{(g)} = \exp\left( \log \pi_\theta(a_t^{(g)} | s_t) - \log \pi_{\theta_{\text{old}}}(a_t^{(g)} | s_t) \right)$ and $\epsilon=0.2$.

\subsubsection{Entropy and KL Regularization}
Additional terms encourage exploration and prevent collapse:
\begin{equation}
    \mathcal{L} = \mathcal{L}_{\text{GRPO}} - \beta_{\text{ent}} \mathbb{E}[\mathcal{H}(\pi_\theta(\cdot|s))] + \beta_{\text{KL}} \text{KL}(\pi_\theta \| \pi_{\text{ref}}),
\end{equation}
where $\pi_{\text{ref}}$ is a reference policy (e.g., behavioral cloning from historical data).

\subsubsection{Why GRPO Suits Microstructure Execution}
GRPO's group evaluation mirrors a trader's mental simulation of alternative execution rates under identical market conditions. The group standard deviation $\sigma_t$ reflects microstructure uncertainty—high volatility or liquidity fragmentation leads to more dispersed rewards, automatically scaling advantages. This provides intrinsic robustness to market noise.

\section{Experiments}
\label{sec:experiments}

\textit{Note: Full experimental results will be reported in the camera-ready version after completion of empirical training. This section describes the experimental setup, baselines, metrics, and planned ablations.}

\subsection{Experimental Setup}

\textbf{Dataset}: We use the FI-2010 limit-order-book dataset \citep{nystrom2014high}, which contains 10 levels of bid/ask prices and volumes for 5 stocks over 10 days. We split data chronologically into training (70\%), validation (15\%), and test (15\%).

\textbf{Execution parameters}: Initial volume $V \sim \mathcal{U}(500, 2000)$ shares, maximum steps $T=100$, impact coefficients $\eta=0.01$, $\psi=0.001$, risk aversion $\alpha=0.1$.

\textbf{Policy architecture}: Encoder MLP with 3 layers of 256 units, conditional RealNVP with 8 coupling layers (hidden dim 128), residual MLP with 2 layers of 128 units.

\textbf{GRPO hyperparameters}: Group size $G=8$, clip range $\epsilon=0.2$, entropy coefficient $\beta_{\text{ent}}=0.01$, KL coefficient $\beta_{\text{KL}}=0.0$ (unless using reference policy).

\subsection{Baselines}
We compare against the following baselines:

\begin{itemize}
    \item \textbf{TWAP}: Time-weighted average price, uniform execution over time.
    \item \textbf{VWAP}: Volume-weighted average price, execution proportional to historical volume.
    \item \textbf{Almgren-Chriss}: Analytical solution with linear impact.
    \item \textbf{Behavioral Cloning (BC)}: Supervised learning on historical executions.
    \item \textbf{Flow-only}: Policy using only the normalizing-flow manifold (no residual).
    \item \textbf{PPO}: Standard Proximal Policy Optimization with value network.
    \item \textbf{CQL}: Conservative Q-Learning \citep{kumar2020conservative}.
    \item \textbf{IQL}: Implicit Q-Learning \citep{kostrikov2022offline}.
\end{itemize}

\subsection{Metrics}
Performance is evaluated using:

\begin{itemize}
    \item \textbf{Implementation shortfall}: Total cost relative to arrival price.
    \item \textbf{Shortfall variance}: Risk measured by cost variability.
    \item \textbf{Slippage}: Difference from mid-price execution.
    \item \textbf{Market impact}: Temporary and permanent impact components.
    \item \textbf{Feasibility score}: Log probability of actions under historical flow model.
\end{itemize}

\subsection{Ablation Studies}
We plan to ablate the following components:

\begin{itemize}
    \item \textbf{No residual module}: Remove the residual adaptation.
    \item \textbf{No flow manifold}: Replace flow with Gaussian policy.
    \item \textbf{Smaller group sizes}: Test $G \in \{2,4,8,16\}$.
    \item \textbf{With critic}: Add a value network for comparison.
    \item \textbf{Different advantage normalization}: Compare group-relative vs. batch-relative.
\end{itemize}

\subsection{Expected Results}
We hypothesize that:

\begin{enumerate}
    \item MACE-RL will outperform traditional benchmarks (TWAP, VWAP, Almgren-Chriss) by adapting to microstructure.
    \item GRPO will show greater training stability than critic-based offline RL methods (CQL, IQL, PPO).
    \item The flow manifold will improve action feasibility and reduce out-of-distribution actions.
    \item The residual module will provide adaptive gains during liquidity shocks.
\end{enumerate}

\textit{Empirical validation of these hypotheses is ongoing.}

\section{Conclusion and Future Work}

We presented MACE-RL, a microstructure-aware conservative execution framework using DeepSeek GRPO. By combining normalizing-flow action manifolds, residual adaptation, and critic-free group policy optimization, our method learns feasible execution strategies directly from historical LOB data. The elimination of value networks addresses a key source of instability in offline RL for finance.

\textbf{Training stability} remains a practical concern; GRPO's group sampling introduces computational overhead and variance that must be managed through careful hyperparameter tuning. Future work could explore adaptive group sizes or hierarchical sampling strategies.

\textbf{Future directions} include:

\begin{itemize}
    \item \textbf{Better impact models}: Incorporating non-linear or stochastic impact functions.
    \item \textbf{Calibrating flows}: Regularizing the flow manifold to ensure economic plausibility.
    \item \textbf{Multi-asset execution}: Extending to portfolio execution with cross-impact.
    \item \textbf{Online adaptation}: Fine-tuning policies with real-time data.
    \item \textbf{Theoretical analysis}: Formal guarantees on GRPO convergence in financial settings.
\end{itemize}

MACE-RL provides a principled foundation for offline RL in optimal execution, bridging microstructure modeling with stable policy optimization.

\bibliographystyle{plain}
\bibliography{references}

\end{document}